\documentclass[
   12pt,
   a4paper,
   twocolumn
]{article}

% Codificação e fontes
\usepackage[utf8]{inputenc}
\usepackage[T1]{fontenc}
\usepackage{lmodern}

% Microtipografia
\usepackage{microtype}

% Idioma
\usepackage[brazil]{babel}

% Margens compactas para caber em 2 páginas
\usepackage[top=2cm, bottom=2cm, left=1.8cm, right=1.8cm]{geometry}

% Figuras
\usepackage{graphicx}
\graphicspath{{sections/figs/}}

% Matemática
\usepackage{amsmath, amssymb}

% Listas compactas
\usepackage{enumitem}
\setlist{nosep, leftmargin=*}

% Espaçamento entre parágrafos e seções compactos
\usepackage{titlesec}
\titlespacing*{\section}{0pt}{6pt}{3pt}
\titlespacing*{\subsection}{0pt}{4pt}{2pt}
\titleformat{\section}{\normalfont\bfseries\normalsize}{\thesection.}{0.5em}{}
\titleformat{\subsection}{\normalfont\bfseries\small}{\thesubsection.}{0.5em}{}

% Sem indentação, com espaço entre parágrafos
\setlength{\parindent}{0pt}
\setlength{\parskip}{4pt}

% Citações
\usepackage[numbers, sort&compress]{natbib}
\bibliographystyle{unsrtnat}

% Hiperlinks
\usepackage{hyperref}
\hypersetup{
  pdftitle={Detecção Automática de Bolhas de Plasma Equatoriais em Imagens de Luminescência Atmosférica Usando Técnicas de Aprendizado de Máquina},
  pdfauthor={Alberto Ivo da Costa Vieira},
  colorlinks=true,
  linkcolor=black,
  citecolor=black,
  urlcolor=blue
}

\begin{document}

% ===== TÍTULO E AUTOR =====
\twocolumn[{%
\begin{center}
{\large\bfseries DETECÇÃO AUTOMÁTICA DE BOLHAS DE PLASMA EQUATORIAIS EM IMAGENS DE LUMINESCÊNCIA ATMOSFÉRICA USANDO TÉCNICAS DE APRENDIZADO DE MÁQUINA\par}
\vspace{6pt}
{\normalsize Alberto Ivo da Costa Vieira\par}
\vspace{12pt}
\end{center}
}]

% ===== SEÇÕES =====
\section{Introdução e Fundamentação Teórica}\label{sec:introduction}

Bolhas de Plasma Equatoriais são grandes regiões de baixa densidade na ionosfera que ocorrem após o pôr do sol. Elas interferem gravemente nos sinais de rádio, prejudicando sistemas de navegação por satélite e comunicações \cite{zhong2025automatic, okoh2025bootstrapping}. Por isso, o problema do monitoramento contínuo desse fenômeno é fundamental para a segurança tecnológica.

A fundamentação teórica deste projeto apoia-se em estudos recentes de clima espacial e visão computacional. Tradicionalmente, a detecção do fenômeno baseia-se na análise manual de imagens do céu noturno capturadas por câmeras especiais. Contudo, esse processo exige muito tempo e esforço, tornando-se inviável diante da enorme quantidade de dados atuais \cite{githio2024machine}. Embora técnicas modernas de inteligência artificial venham sendo testadas, a literatura demonstra que o uso direto de imagens brutas exige alto poder de processamento e frequentemente confunde os algoritmos com ruídos visuais do ambiente, como nuvens e a luz da lua.

Justifica-se, portanto, a necessidade de propor o desenvolvimento de um método de inteligência artificial delimitado e específico para imagens noturnas. A proposta prepara as fotografias para remover as luzes de fundo, reduz a complexidade dos dados de forma inteligente e combina algoritmos avançados para identificar as bolhas com precisão e velocidade, viabilizando o monitoramento em tempo real.
\section{Objetivos}\label{sec:objectives}

O objetivo central da pesquisa é criar e avaliar modelos de inteligência artificial que identifiquem automaticamente as bolhas em imagens do céu noturno do programa EMBRACE do INPE, aprimorando o monitoramento do clima espacial brasileiro.

Os passos práticos incluem:
\begin{enumerate}
   \item Organizar um conjunto de imagens do INPE, contendo noites com e sem bolhas sob diferentes condições climáticas.
   \item Preparar as imagens removendo luzes indesejadas e destacando as estruturas de interesse.
   \item Diminuir o tamanho dos arquivos de imagem matematicamente, mantendo o formato e a localização espacial das informações.
   \item Treinar um comitê de inteligência artificial que une três algoritmos diferentes, utilizando estratégias para equilibrar os dados de noites normais e noites com bolhas.
   \item Validar as decisões da máquina para confirmar se ela avalia o formato correto do fenômeno e testar sua estabilidade durante monitoramentos contínuos de longos períodos.
\end{enumerate}
\section{Metodologia}\label{sec:methodology}

A metodologia une a física da atmosfera e a inteligência artificial para lidar com as imagens capturadas pelo INPE, organizada em fases de pesquisa plenamente factíveis de serem executadas no período de 24 meses do Mestrado.

O primeiro passo é o preparo. Devido à iluminação irregular do céu noturno, aplicaremos filtros para suavizar e subtrair o brilho difuso. Em seguida, melhoraremos o contraste para que os contornos das bolhas se destaquem contra o fundo escuro.

Como a maioria das noites não apresenta bolhas, usaremos estratégias para equilibrar os dados, multiplicando imagens e gerando novos exemplos para impedir que o algoritmo fique viciado em prever o céu limpo. Para reduzir o peso dos arquivos sem perder a qualidade espacial da fotografia, aplicaremos a Análise de Componentes Principais Bidimensional, preservando características e acelerando o processamento.

A classificação final será validada por um comitê de inteligência artificial unindo três modelos em equipe. Ferramentas de transparência garantirão que a máquina de fato analisa o formato das bolhas.

\section{Resultados Preliminares}\label{sec:results}

Os testes iniciais alcançaram uma taxa de acerto próxima a cem por cento no reconhecimento das bolhas, distinguindo com grande precisão as noites normais das noites afetadas pelo fenômeno climático.

Ao utilizarmos as ferramentas de explicação visual, constatamos que a máquina focou sua atenção consistentemente na parte superior do céu, exatamente a região onde as bolhas costumam ascender e se manifestar na realidade. Isso comprova que o modelo aprendeu o padrão físico autêntico.

Além disso, simulamos o funcionamento do modelo de forma ininterrupta durante uma noite inteira. O sistema se comportou de maneira impecável ao processar centenas de fotos limpas consecutivas sem gerar nenhum falso positivo, reforçando sua vocação para atuar no monitoramento constante e autônomo.

\section{Considerações Finais}\label{sec:conclusion}

A pesquisa inicial demonstrou que tratar as características visuais do céu e simplificar os dados antes do aprendizado resulta em um sistema muito mais seguro e transparente do que simplesmente utilizar imagens brutas. O método foi preciso e conseguiu explicar logicamente suas escolhas.

Durante o mestrado, testaremos o programa com fotos de outras regiões do país e sob diferentes comportamentos da atividade solar. Também investigaremos como combinar essas imagens com dados capturados por radares e sistemas GPS, oferecendo uma solução completa e de alto nível para os centros de monitoramento diário.


% ===== REFERÊNCIAS =====
\vspace{4pt}
{\small
\bibliography{sections/biblio}
}

\end{document}
